%! Graphing Exponential Functions and Transformations — Unit 6, Lesson 6.2 — Group Activity
\documentclass[10pt]{article}
\usepackage{algebra2-article}
\usepackage{algebra2-boxes}
\newcommand{\TallMath}[1]{$\displaystyle #1\rule[-1.4em]{0pt}{3.2em}$}
\newcommand{\lgrid}[4]{%
  \draw[step=1, linegray!40, very thin] (#1,#3) grid (#2,#4);
  \draw[->, semithick, charcoal] (#1,0)--(#2,0) node[below right, font=\scriptsize]{$x$};
  \draw[->, semithick, charcoal] (0,#3)--(0,#4) node[left, font=\scriptsize]{$y$};}
% g(x) = a*b^(x-h) + k plotted as a*exp((ln b)(x-h))+k.
% #1=a  #2=ln(b)  #3=h  #4=k  #5=xmin  #6=xmax
% ln2 = 0.693147   ln3 = 1.098612   ln(1/2) = -0.693147
\newcommand{\expgraph}[6]{\draw[very thick, forest, domain=#5:#6, samples=140, smooth]
  plot (\x,{(#1)*exp((#2)*((\x)-(#3)))+(#4)});}
\newcommand{\hasym}[3]{\draw[navy, dashed, semithick] (#1,#3)--(#2,#3);}

\begin{document}

\pageheader{Unit 6, Lesson 6.2}{Group Activity}
\namepartnerperiod

\vspace{0.08in}

\begin{headlinebox}{forestbg}
\small\textbf{Move It --- and Say What Moved.} Every problem below starts from a parent $y=b^{x}$ and
does something to it. For each one, your partner should be able to stop you and ask the two questions
that decide everything: \textbf{did the asymptote move?} and \textbf{what is the range now?} If a
transformation did not move the dashed line, it did not change the range either.
\end{headlinebox}

\vspace{0.06in}

\begin{tcolorbox}[colback=white, colframe=black!40, title=\textbf{Tier R --- Remediate},
    fonttitle=\bfseries, arc=1.5mm, left=3mm, right=3mm, top=2mm, bottom=2mm, breakable]
\textbf{Part 1 --- What moved?} Every row starts from the parent $y=3^{x}$, whose asymptote is $y=0$
and whose range is $(0,\infty)$.
\begin{center}
\renewcommand{\arraystretch}{1.6}
\begin{tabularx}{\linewidth}{@{}l X l X@{}}
\hline
\textbf{Function} & \textbf{Move(s) from the parent} & \textbf{Asymptote} & \textbf{Range} \\
\hline
$y=3^{x}+4$ & \blank{4.0cm} & \blank{1.6cm} & \blank{2.4cm} \\
$y=3^{\,x-2}$ & \blank{4.0cm} & \blank{1.6cm} & \blank{2.4cm} \\
$y=3^{x}-5$ & \blank{4.0cm} & \blank{1.6cm} & \blank{2.4cm} \\
$y=-3^{x}$ & \blank{4.0cm} & \blank{1.6cm} & \blank{2.4cm} \\
\hline
\end{tabularx}
\end{center}
Two of those four rows have the same asymptote as the parent. Which two, and what do they have in
common? \blank{7.5cm}

\vspace{6pt}
\textbf{Part 2 --- One function, all the way through.} Let $g(x)=2^{x}+3$.
\begin{enumerate}[label=(\alph*), leftmargin=1.9em, itemsep=6pt]
  \item Parent function: $y=$ \blank{1.6cm} \quad Move: \blank{3.0cm}
  \item Horizontal asymptote: $y=$ \blank{1.4cm}
  \item Fill in the table. (Remember $2^{-1}=\frac12$ and $2^{-2}=\frac14$.)
    \begin{center}
    \renewcommand{\arraystretch}{1.35}
    \begin{tabular}{|c|c|c|c|c|c|}
    \hline
    $x$ & $-2$ & $-1$ & $0$ & $1$ & $2$ \\ \hline
    $g(x)$ & \blank{1.3cm} & \blank{1.3cm} & \blank{1.3cm} & \blank{1.3cm} & \blank{1.3cm} \\ \hline
    \end{tabular}
    \end{center}
  \item $y$-intercept: \blank{1.8cm} \quad Range: \blank{2.4cm}
  \item Every output in your table is bigger than $3$. Will \emph{any} input ever give an output of
        exactly $3$? Why not? \blank{5.6cm}
  \item Does this graph have an $x$-intercept? Explain using the asymptote.
        \blank{5.6cm}
\end{enumerate}
\end{tcolorbox}

\begin{tcolorbox}[colback=white, colframe=black!40, title=\textbf{Tier A --- Approaching Proficiency},
    fonttitle=\bfseries, arc=1.5mm, left=3mm, right=3mm, top=2mm, bottom=2mm, breakable]
\textbf{Part 1 --- Match, then read.} Each graph below is a transformation of $y=2^{x}$. The dashed
line is the horizontal asymptote.
\begin{center}
\begin{tikzpicture}[scale=0.28]
  % P: y = 2^x + 2
  \begin{scope}
    \lgrid{-2}{5}{-4}{8}
    \begin{scope}\clip (-2,-4) rectangle (5,8); \expgraph{1}{0.693147}{0}{2}{-2}{5} \end{scope}
    \hasym{-2}{5}{2}
    \node[charcoal, font=\small] at (1.5,-5.6) {\textbf{P}};
  \end{scope}
  % Q: y = 2^x - 2
  \begin{scope}[xshift=9.5cm]
    \lgrid{-2}{5}{-4}{8}
    \begin{scope}\clip (-2,-4) rectangle (5,8); \expgraph{1}{0.693147}{0}{-2}{-2}{5} \end{scope}
    \hasym{-2}{5}{-2}
    \node[charcoal, font=\small] at (1.5,-5.6) {\textbf{Q}};
  \end{scope}
  % R: y = 2^(x-2)
  \begin{scope}[xshift=19cm]
    \lgrid{-2}{5}{-4}{8}
    \begin{scope}\clip (-2,-4) rectangle (5,8); \expgraph{1}{0.693147}{2}{0}{-2}{5} \end{scope}
    \hasym{-2}{5}{0}
    \node[charcoal, font=\small] at (1.5,-5.6) {\textbf{R}};
  \end{scope}
  % S: y = 2^(-x)
  \begin{scope}[xshift=28.5cm]
    \lgrid{-2}{5}{-4}{8}
    \begin{scope}\clip (-2,-4) rectangle (5,8); \expgraph{1}{-0.693147}{0}{0}{-2}{5} \end{scope}
    \hasym{-2}{5}{0}
    \node[charcoal, font=\small] at (1.5,-5.6) {\textbf{S}};
  \end{scope}
\end{tikzpicture}
\end{center}
\vspace{2pt}
Write the letter of the matching graph, then fill in the last two columns from the \emph{equation}.
\begin{center}
\renewcommand{\arraystretch}{1.55}
\begin{tabularx}{\linewidth}{@{}l c X X@{}}
\hline
\textbf{Equation} & \textbf{Graph} & \textbf{Asymptote} & \textbf{Range} \\
\hline
$y=2^{x}+2$ & \blank{1.0cm} & \blank{2.2cm} & \blank{2.6cm} \\
$y=2^{x}-2$ & \blank{1.0cm} & \blank{2.2cm} & \blank{2.6cm} \\
$y=2^{\,x-2}$ & \blank{1.0cm} & \blank{2.2cm} & \blank{2.6cm} \\
$y=2^{-x}$ & \blank{1.0cm} & \blank{2.2cm} & \blank{2.6cm} \\
\hline
\end{tabularx}
\end{center}
Exactly one of those four graphs has an $x$-intercept. Which one, and why is it the only one?
\blank{7.5cm}

\vspace{6pt}
\textbf{Part 2 --- Backwards: from the graph to the equation.} The curve has the form
$y=a\,b^{x}+k$.
\begin{center}
\begin{minipage}[c]{0.42\linewidth}
\centering
\begin{tikzpicture}[scale=0.40]
  \lgrid{-3}{4}{-3}{11}
  \begin{scope}\clip (-3,-3) rectangle (4,11); \expgraph{3}{0.693147}{0}{-1}{-3}{4} \end{scope}
  \hasym{-3}{4}{-1}
  \node[navy, font=\scriptsize, left] at (-1.2,-1.7) {$y=-1$};
  \fill[charcoal] (0,2) circle (3pt) node[left, font=\scriptsize]{$(0,2)$};
  \fill[charcoal] (1,5) circle (3pt) node[right, font=\scriptsize]{$(1,5)$};
\end{tikzpicture}
\end{minipage}
\begin{minipage}[c]{0.54\linewidth}
Step 1 --- the dashed line gives $k=$ \blank{1.2cm} \\[6pt]
Step 2 --- the $y$-intercept is $a+k$, so $a=$ \blank{1.2cm} \\[6pt]
Step 3 --- use $(1,5)$: \; $5=$ \blank{1.0cm}$\cdot b+$ \blank{1.0cm}, so $b=$ \blank{1.2cm} \\[6pt]
Equation: $y=$ \blank{3.0cm} \\[6pt]
Range: \blank{2.4cm}
\end{minipage}
\end{center}

\vspace{4pt}
\textbf{Justify.} Nadia says $y=-5^{x}$ and $y=5^{-x}$ are ``the same thing, just written
differently --- they both have a negative.'' She is wrong. For \emph{each} function, say which axis
it reflects over, give its range, and say whether it is growth or decay. Then state in one sentence
what Nadia is actually confusing. \writelines{4}
\end{tcolorbox}

\begin{tcolorbox}[colback=white, colframe=black!40, title=\textbf{Tier E --- Extension},
    fonttitle=\bfseries, arc=1.5mm, left=3mm, right=3mm, top=2mm, bottom=2mm, breakable]
\textbf{Part 1 --- A horizontal shift in disguise.} Exponentials do something no other family in this
course does. Use the exponent rules from Unit 5.
\begin{enumerate}[label=(\alph*), leftmargin=1.9em, itemsep=6pt]
  \item Rewrite $2^{\,x-3}$ as a product: \; $2^{\,x-3}=2^{x}\cdot 2^{\,\square}$ where $\square=$
        \blank{1.0cm}, \; and $2^{\,\square}=$ \blank{1.2cm}. \\[3pt]
        So $2^{\,x-3}=$ \blank{1.4cm} $\cdot\, 2^{x}$.
  \item Check it in a table.
    \begin{center}
    \renewcommand{\arraystretch}{1.35}
    \begin{tabular}{|c|c|c|c|c|}
    \hline
    $x$ & $3$ & $4$ & $5$ & $6$ \\ \hline
    $2^{\,x-3}$ & \blank{1.2cm} & \blank{1.2cm} & \blank{1.2cm} & \blank{1.2cm} \\ \hline
    $\frac18\cdot 2^{x}$ & \blank{1.2cm} & \blank{1.2cm} & \blank{1.2cm} & \blank{1.2cm} \\ \hline
    \end{tabular}
    \end{center}
  \item So for an exponential, shifting \emph{right} $3$ produces the same graph as a vertical
        \blank{2.6cm} by a factor of $\frac18$. Does either move change the asymptote?
        \blank{2.0cm}
  \item Try the same trick on the parabola: is $(x-3)^{2}$ equal to some fixed number times $x^{2}$?
        Test it at $x=4$ and $x=5$ and explain what goes wrong. \writelines{2}
\end{enumerate}

\vspace{5pt}
\textbf{Part 2 --- All four moves at once.} Let $g(x)=-2(3)^{\,x-1}+5$.
\begin{center}
\renewcommand{\arraystretch}{1.5}
\begin{tabularx}{\linewidth}{@{}l X@{}}
\hline
Parent function & \blank{3.0cm} \\
The moves, in words & \blank{9.0cm} \\
Horizontal asymptote & \blank{3.0cm} \\
Range & \blank{3.0cm} \\
$y$-intercept ($g(0)$, exact) & \blank{3.0cm} \\
\hline
\end{tabularx}
\end{center}
Does $g$ have an $x$-intercept? Answer using the signs of $a$ and $k$ --- not by graphing.
\blank{7.0cm}

\vspace{5pt}
\textbf{Part 3 --- When the asymptote is a real thing.} A cup of coffee is poured at
$200^{\circ}$F into a room held at $68^{\circ}$F. Its temperature after $t$ minutes is
\[
  T(t) = 68 + 132(0.85)^{t}.
\]
\begin{center}
\renewcommand{\arraystretch}{1.3}
\begin{tabular}{|c|c|c|c|c|c|c|}
\hline
$t$ (min) & $0$ & $4$ & $8$ & $12$ & $16$ & $20$ \\ \hline
$T(t)$ ($^{\circ}$F) & $200$ & $136.9$ & $104.0$ & \blank{1.2cm} & $77.8$ & $73.1$ \\ \hline
\end{tabular}
\end{center}
\begin{enumerate}[label=(\alph*), leftmargin=1.9em, itemsep=6pt]
  \item Identify $a$, $b$, and $k$: \; $a=$ \blank{1.4cm} \quad $b=$ \blank{1.4cm} \quad
        $k=$ \blank{1.4cm}
  \item Fill the missing table entry by evaluating $T(12)$. \blank{3.0cm}
  \item The horizontal asymptote is $y=$ \blank{1.4cm}. What does that number \emph{mean} about
        this coffee? \blank{5.0cm}
  \item Algebraically the range of $T$ is $(68,\infty)$. But the temperatures this coffee actually
        takes are only $(68,200]$. Explain the difference. \writelines{2}
  \item At what time is the coffee exactly $68^{\circ}$F? \blank{5.0cm}
  \item Between which two table times does the coffee pass $100^{\circ}$F? \blank{2.6cm} \\[3pt]
        To find that moment \emph{exactly} you would have to solve $132(0.85)^{t}=32$. Try rewriting
        both sides with a common base --- what goes wrong? \writelines{2}
\end{enumerate}
\end{tcolorbox}

\end{document}
